% !PRINTED 726
\weberpage{726}

\begin{center}
\S.\,199.

\gesperrt{Der Classenzahlfactor} $B$.
\end{center}

Auf ganz andere Weise muss der Classenzahlfactor $B$ berechnet werden. Hier ist, wenn $\Theta$ wieder eine Primitivwurzel
der Gleichung \S.\,198, (3) ist ($\Theta^{\nu'} + 1 = 0$), und
\weq{1}{t \equiv (-1)^\alpha\,5^\beta \pmod{m},}
\weq{2}{\chi_1(t) = \Theta^\beta}
zu setzen, und die Formel \S.\,197, (15) ergiebt:
\weq{3}{X_1 = -\frac{2}{m}\,(1, \Theta^{-1}, r)\,\overset{t}{\Sigma}\,\Theta^\beta\,\log\,\mathrm{tg}\,\frac{t\pi}{m}.}

Nun ist für jedes ungerade $t$, wenn [mit Rücksicht auf (1)]
\[
r = e^{\frac{2\pi i}{m}}, \quad r^\nu = i, \quad r^{\nu t} = (-1)^\alpha i
\]
gesetzt wird,
\weq{4}{\mathrm{tg}\,\frac{t\pi}{m} = (-1)^\alpha\,r^{t\nu}\,\frac{1 - r^t}{1 + r^t},}
und dies ist, wie wir schon im \S.\,178 gesehen haben, eine Einheit des reellen Körpers $H_m$. Man erhält die conjugirten Werthe
dieser Einheit in Bezug auf diesen Körper, wenn man $t \equiv 5^\beta \pmod{m}$ setzt, und $\beta$ ein volles Restsystem nach dem Modul $\mu$
durchlaufen lässt. Wir setzen demnach:
\weq{5}{\tau_\beta = \mathrm{tg}\,\frac{5^\beta\pi}{m}.}

Nach \S.\,198, (7) ist $5^{\beta+\nu'} \equiv 5^\beta + \mu \pmod{m}$, und daraus
ergiebt sich nach einer bekannten trigonometrischen Formel:
\weq{6}{\tau_{\beta+\nu'} = \frac{-1}{\tau_\beta}.}

Setzen wir also nun
\weq{7}{\lambda_\beta = \tfrac{1}{2}\,\log\,\tau_\beta^2,}
so ist nach (6)
\weq{8}{\lambda_{\beta+\nu'} = -\lambda_\beta,}
und demnach
\weq{9}{\Theta^{\beta+\nu'}\,\lambda_{\beta+\nu'} = \Theta^\beta\,\lambda_\beta;}
unter den in der Summe (3) vorkommenden Werthen von $\beta$
sind, wie wir schon im vorigen Paragraphen gesehen haben, nicht
zwei nach dem Modul $\nu'$ congruent, und wenn wir daher mittelst
%CONT>>
